% from ega2-1-fr.tex
\oldpage[II]{118}

$B$-algèbre~; soit $B'$ la fermeture intégrale de $B$ dans $A$. Tout revient
à voir que pour tout $x\in X$, un élément de $A_x$, entier sur $B_x$,
appartient nécessairement à $B'_x$, ce qui résulte du fait que, pour un anneau
commutatif $C$, les opérations de fermeture intégrale dans une $C$-algèbre et
de passage à un anneau de fractions (pour une partie multiplicative de $C$)
commutent (\cite[t.~I, p.~261 et 257]{I-13}).

Le $X$-schéma $X'=\operatorname{Spec}(\mathcal{O}'_X)$ est alors appelé la
\emph{fermeture intégrale de $X$ relativement à $\mathcal{A}$} (ou
\emph{relativement à $\operatorname{Spec}(\mathcal{A})$})~; il est clair que
$X'$ est \emph{entier} sur $X$ (\hyperref[II.6.1.2-fr]{6.1.2}).

On déduit aussitôt de (\hyperref[II.6.3.4-fr]{6.3.4}) que si $f:X'\to X$ est
le morphisme structural, alors, pour tout ouvert $U$ de $X$, $f^{-1}(U)$ est
la \emph{fermeture intégrale du préschéma induit par $X$ sur $U$, relativement
à $\mathcal{A}|U$}.

\begin{env}[6.3.5]
\label{II.6.3.5-fr}
Soient $X$, $Y$ deux préschémas, $f:X\to Y$ un morphisme, $\mathcal{A}$
(resp. $\mathcal{B}$) une $\mathcal{O}_X$-Algèbre (resp. une
$\mathcal{O}_Y$-Algèbre) quasi-cohérente, et soit
$u:\mathcal{B}\to\mathcal{A}$ un $f$-morphisme. On a vu
(\hyperref[II.6.3.3-fr]{6.3.3}) qu'on en déduit un $f$-morphisme
$u':\mathcal{O}'_Y\to\mathcal{O}'_X$, où $\mathcal{O}'_X$ (resp.
$\mathcal{O}'_Y$) est la fermeture intégrale de $\mathcal{O}_X$ (resp.
$\mathcal{O}_Y$) dans $\mathcal{A}$ (resp. $\mathcal{B}$). Par suite, si
$X'$ (resp. $Y'$) est la fermeture intégrale de $X$ (resp. $Y$) relativement
à $\mathcal{A}$ (resp. $\mathcal{B}$), on déduit canoniquement de $u$ un
morphisme $f'=\operatorname{Spec}(u'):X'\to Y'$
(\hyperref[II.1.5.6-fr]{1.5.6}) rendant commutatif le diagramme
\[
\label{II.6.3.5.1-fr}
\xymatrix{
X'\ar[r]^{f'}\ar[d]&Y'\ar[d]\\
X\ar[r]_{f}&Y
}
\tag{6.3.5.1}
\]
\end{env}

\begin{env}[6.3.6]
\label{II.6.3.6-fr}
Supposons que $X$ n'ait qu'un nombre \emph{fini} de composantes irréductibles
$X_i$ ($1\leq i\leq r$), de points génériques $\xi_i$, et considérons en
particulier la fermeture intégrale de $X$ relativement à une
$\mathcal{R}(X)$-Algèbre $\mathcal{A}$ quasi-cohérente (en tant que
$\mathcal{O}_X$-Algèbre ou en tant que $\mathcal{R}(X)$-Algèbre, ce qui
revient au même). On sait (\hyperref[I.7.3.5-fr]{I, 7.3.5}) que
$\mathcal{A}$ est composée directe de $r$ $\mathcal{O}_X$-Algèbres
quasi-cohérentes $\mathcal{A}_i$, le support de $\mathcal{A}_i$ étant contenu
dans $X_i$, et le faisceau induit par $\mathcal{A}_i$ sur $X_i$ étant un
faisceau simple dont la fibre $A_i$ est une algèbre sur
$\mathcal{O}_{\xi_i}$. Il est clair alors (\hyperref[II.6.3.4-fr]{6.3.4}) que
la fermeture intégrale $\mathcal{O}'_X$ de $\mathcal{O}_X$ dans
$\mathcal{A}$ est \emph{composée directe} des fermetures intégrales
$\mathcal{O}^{(i)}_X$ de $\mathcal{O}_X$ dans chacune des $\mathcal{A}_i$, et
par suite la fermeture intégrale $X'=\operatorname{Spec}(\mathcal{O}'_X)$ de
$X$ relativement à $\mathcal{A}$ est un $X$-schéma \emph{somme} des
$\operatorname{Spec}(\mathcal{O}^{(i)}_X)=X'_i$ ($1\leq i\leq r$).

Supposons en outre que la $\mathcal{O}_X$-Algèbre $\mathcal{A}$ soit
\emph{réduite}, ou, ce qui revient au même, que chacune des algèbres $A_i$
soit réduite, et puisse par suite être considérée comme une algèbre sur le
\emph{corps} $k(\xi_i)$ (égal au corps des fonctions rationnelles du
sous-préschéma réduit de $X$ ayant $X_i$ pour espace sous-jacent)~; alors
(\hyperref[II.1.3.8-fr]{1.3.8}) chacun des $X'_i$ est un $X$-schéma
\emph{réduit} et $X'$ est aussi la fermeture intégrale de $X_{\mathrm{red}}$.
Supposons de plus que chacune des algèbres $A_i$ soit \emph{composée directe
d'un nombre fini de corps} $K_{ij}$ ($1\leq j\leq s_i$)~; si
$\mathcal{K}_{ij}$ est la sous-Algèbre de $\mathcal{A}_i$ correspondant à
$K_{ij}$, il est clair que $\mathcal{O}^{(i)}_X$ est \emph{composée directe}
des fermetures intégrales $\mathcal{O}^{(ij)}_X$ de $\mathcal{O}_X$ dans
chacune des $\mathcal{K}_{ij}$. Par suite, $X'_i$ est alors \emph{somme} des
$X$-schémas $X'_{ij}=\operatorname{Spec}(\mathcal{O}^{(ij)}_X)$
($1\leq j\leq s_i$). En outre, sous ces hypothèses et avec ces notations~:
\end{env}

